\begin{titlepage}
\centering
\vspace*{2cm}
\includegraphics[width=5cm]{logo}\par
\vspace{1.5cm}
{\sffamily\bfseries\fontsize{38}{44}\selectfont\color{tinta} Manual de\\[0.2em] Lázaro Linux\par}
\vspace{0.8cm}
{\sffamily\Large\color{pulso} Tu hardware aún tiene pulso\par}
\vspace{1.5cm}
{\large Guía para instalar y usar Lázaro Linux en el día a día,\\
sin conocimientos técnicos\par}
\vfill
{\sffamily Versión \lazaroversion\ \textperiodcentered\ \lazarofecha\par}
\end{titlepage}

\thispagestyle{empty}
\vspace*{\fill}
\noindent{\small
\textbf{Manual de Lázaro Linux}, versión \lazaroversion.\\
Proyecto Lázaro Linux, \url{https://github.com/Entropio1984/lazaro-linux}.

\medskip
\noindent Este manual es libre: puedes copiarlo, imprimirlo, compartirlo y
mejorarlo, siempre que digas de dónde viene y compartas tus cambios con la
misma licencia, la Creative Commons Atribución-CompartirIgual 4.0
Internacional (CC BY-SA 4.0):
\url{https://creativecommons.org/licenses/by-sa/4.0/deed.es}.

\medskip
\noindent Los nombres de otros programas y marcas que aparecen aquí son de sus
dueños y se mencionan solo para explicar cómo usarlos.

\medskip
\noindent ¿Encontraste un error o falta algo? Avísanos en
\url{https://github.com/Entropio1984/lazaro-linux/issues}.\par}
\cleardoublepage
